%% Paper 1 -- Insurance: Mathematics and Economics (IME)
%% Elsevier elsarticle template. Compile: pdflatex; bibtex; pdflatex x2.
\documentclass[review,3p,times]{elsarticle}

\usepackage{amsmath,amssymb,amsthm}
\usepackage{booktabs}
\usepackage{graphicx}
\graphicspath{{../output/figures/}}
\usepackage[hidelinks]{hyperref}

\newtheorem{proposition}{Proposition}
\newcommand{\dd}{\mathrm{d}}
\newcommand{\E}{\mathbb{E}}

\journal{Insurance: Mathematics and Economics}

\begin{document}

\begin{frontmatter}

\title{A Bayesian quadratic-hazard model for biomarker-driven mortality
and morbidity risk: optimal health levels with uncertainty}

\author[byu]{R. Richardson}
\author[bio]{[Biologist co-author]} % Paper 2 lead; advisory here if appropriate
\address[byu]{Department of Statistics, Brigham Young University, Provo, UT, USA}
\address[bio]{[Department], Brigham Young University, Provo, UT, USA}

\begin{abstract}
Insurers increasingly underwrite and price on repeatedly measured biological
markers---body mass index, blood pressure, glucose, lipids---yet standard
actuarial hazard models impose monotone or multiplicative marker effects that
cannot capture the U- and J-shaped mortality and morbidity risks pervasive in
epidemiology. We develop a Bayesian quadratic-hazard stochastic process model
that jointly describes (i) the irregular longitudinal dynamics of a biomarker
through a mean-reverting (Ornstein--Uhlenbeck) diffusion and (ii) a quadratic
hazard whose minimum identifies an age-varying risk-minimizing level
$f_0(t)$---a ``physiological optimum''---with curvature $Q(t)$ and mean-reversion
$a(t)$ interpretable as stress-resistance and adaptive capacity. Estimation
exploits the analytically marginalized likelihood, a Kalman-type moment recursion
that integrates the quadratic hazard against the Gaussian latent state in closed
form; we embed this in a Bayesian framework to obtain full posterior uncertainty.
The treatment delivers (a) credible bands on the optimal trajectory $f_0(t)$ and
the latent risk components, and (b) posterior-predictive, dynamically updated
survival and hazard forecasts with coherent uncertainty---quantities of direct
relevance to biomarker-based underwriting, wellness-program design, and the
uncertainty margins required under modern solvency and reporting regimes. We show
that priors are essential for identifiability rather than cosmetic: on
monotone-risk markers the maximum-likelihood estimator collapses the quadratic
term, while weakly-informative priors recover a stable, interpretable fit. The
methods are implemented in \texttt{R}, validated by simulation, and illustrated
on publicly available longitudinal survival data.
\end{abstract}

\begin{keyword}
Longevity and mortality risk \sep Biomarkers \sep Joint models \sep
Bayesian inference \sep Uncertainty quantification \sep Underwriting
\MSC[2020] 91G05 \sep 62F15 \sep 62N02
\end{keyword}

\end{frontmatter}

%====================================================================
\section{Introduction}\label{sec:intro}
% Biomarker/accelerated underwriting and health-insurance context.
% Standard actuarial hazard models (PH-GLM, aggregate mortality) assume
% monotone/multiplicative marker effects; epidemiology shows U-/J-shaped risk.
% Longitudinal markers are irregular/missing; their dynamics carry risk info.
% The need: joint biomarker--event modeling WITH uncertainty for risk margins.
% Existing SPM/quadratic-hazard work is frequentist (NAAJ 2012, 2016); we give
% the Bayesian formulation. Contributions (bullet list).

%====================================================================
\section{The model}\label{sec:model}
The biomarker $Y(t)$ follows a mean-reverting diffusion
\begin{equation}\label{eq:sde}
  \dd Y(t) = a(t)\,\bigl[Y(t)-f_1(t)\bigr]\,\dd t + b(t)\,\dd W(t),
\end{equation}
and the force of mortality (or morbidity onset) is quadratic in the marker,
\begin{equation}\label{eq:hazard}
  \mu(t,Y) = \mu_0(t) + Q(t)\,\bigl[Y(t)-f_0(t)\bigr]^2 .
\end{equation}
% Actuarial interpretation: f0 = risk-minimizing ("optimal") level; Q =
% sensitivity of risk to deviation (stress-resistance/robustness); a =
% recovery rate (adaptive capacity); f1 = drift/allostatic mean; b = volatility.
% State component parameterizations (linear-in-age etc.).

%====================================================================
\section{The marginalized likelihood}\label{sec:lik}
\begin{proposition}\label{prop:recursion}
% Conditional on survival and the observation history, the latent state remains
% Gaussian with mean m(t), variance gamma(t) solving
%   m'  = a(t)(m-f1) - 2 Q(t) gamma (m-f0),
%   g'  = 2 a(t) gamma + b(t)^2 - 2 Q(t) gamma^2,
% and the effective hazard given survival is mu0 + Q[(m-f0)^2 + gamma].
\end{proposition}
% Proof sketch / reference (Woodbury--Manton 1977; Yashin et al. 2007, 2012).
% Observation update with measurement error tau (classic SPM = tau -> 0).
% Likelihood contribution per subject; remark on irregular/missing data.

%====================================================================
\section{Bayesian inference}\label{sec:bayes}
% Priors (weakly-informative) and the identifiability argument: MLE collapses
% the quadratic term on monotone-risk data (aQ -> 0, singular information);
% priors restore a stable, interpretable posterior. Computation: HMC/NUTS on
% the marginal log-posterior (RTMB objective via tmbstan). Model comparison:
% WAIC/PSIS-LOO in place of AIC.

%====================================================================
\section{Actuarial quantities}\label{sec:actuarial}
% (a) Credible bands on f0(t): the optimal metric and its uncertainty.
% (b) Posterior-predictive dynamic survival/hazard via landmarking.
% (c) Survival / risk by deviation-from-optimal, with credible intervals.
% Link each to an actuarial use: underwriting, reserving margin, wellness.

%====================================================================
\section{Simulation study}\label{sec:sim}
% Parameter recovery; coverage/calibration of credible bands; the role of
% priors for identifiability.

%====================================================================
\section{Illustration: longitudinal survival data}\label{sec:pbc}
% PBC (survival::pbcseq): real, irregular, partially-missing data; MAP vs the
% MLE collapse; estimated optimal marker trajectory on the natural scale.
% Figure: pbc_f0_bilirubin.png.

%====================================================================
\section{Discussion}\label{sec:disc}
% Risk margins (Solvency II) / risk adjustment (IFRS 17); wellness incentives;
% limitations (single biomarker, parametric components); extensions
% (multivariate markers, underwriting covariates, forecasting); pointer to the
% CHS case study (companion paper).

%====================================================================
\section*{Acknowledgements}
% Funding; data sources for the public illustration.

\bibliographystyle{elsarticle-harv}
\bibliography{refs}

\end{document}
